% Generated draft: requires mathematical and source review.
\documentclass[11pt]{article}
\usepackage[margin=25mm]{geometry}
\usepackage{fontspec}
\setmainfont{Latin Modern Roman}
\usepackage{amsmath,amssymb,mathtools,unicode-math}
\setmathfont{Latin Modern Math}
\usepackage{enumitem,needspace,xurl,hyperref}
\setlength{\parindent}{0pt}
\setlength{\parskip}{5pt}
\setlength{\emergencystretch}{3em}
\newcommand{\N}{\mathbb{N}}
\newcommand{\Z}{\mathbb{Z}}
\newcommand{\Q}{\mathbb{Q}}
\newcommand{\R}{\mathbb{R}}
\newcommand{\C}{\mathbb{C}}
\newcommand{\F}{\mathbb{F}}
\newcommand{\A}{\mathbb{A}}
\newcommand{\PP}{\mathbb{P}}
\newcommand{\E}{\mathbb{E}}
\newcommand{\eps}{\varepsilon}
\newcommand{\dd}{\,\mathrm{d}}
\begin{document}
\section*{UnsolvedMath ID: 701581}
\textbf{Source problem code:} catalogue entry
\subsection{Definitions and mathematical statement}
Let \(G\) be a finite group, let \(p\) be a prime, and let \(P\leq G\) be a Sylow \(p\)-subgroup. For a finite group \(H\), let \(\operatorname{Irr}(H)\) denote its irreducible ordinary complex characters, and set
\[
\operatorname{Lin}(H)=\{\lambda\in\operatorname{Irr}(H):\lambda(1)=1\}.
\]
For \(\chi\in\operatorname{Irr}(G)\) and \(\lambda\in\operatorname{Lin}(P)\), the multiplicity of \(\lambda\) in the restriction \(\chi|_P\) is
\[
m_\lambda=\langle\chi|_P,\lambda\rangle_P
=\frac{1}{|P|}\sum_{x\in P}\chi(x)\overline{\lambda(x)}.
\]
Define the set of distinct linear constituents by
\[
L_P(\chi)=\{\lambda\in\operatorname{Lin}(P):m_\lambda>0\}.
\]
\textbf{Conjecture (E. Giannelli, Kourovka Notebook Problem 19.33).} For every such \(G,p,P\) and every \(\chi\in\operatorname{Irr}(G)\),
\[
p\mid\chi(1)\quad\text{and}\quad L_P(\chi)\ne\varnothing
\quad\Longrightarrow\quad |L_P(\chi)|\geq p.
\]
Thus the conclusion counts distinct linear characters, regardless of their multiplicities.

The supplied source reports that the conjecture has been verified for symmetric groups, alternating groups, \(p\)-solvable groups, and sporadic simple groups. Here a finite group is \(p\)-solvable if it admits a normal series whose factors are each either of prime-power order \(p^a\) or of order relatively prime to \(p\). These reported verifications are source qualifications; their proofs and specific supporting references were not supplied, and they have not been independently verified here. The registry labels the general conjecture as open.
\subsection{Short English statement}
Giannelli's conjecture states that if an irreducible ordinary character of a finite group has degree divisible by a prime \(p\), then its restriction to a Sylow \(p\)-subgroup has either no linear constituents or at least \(p\) distinct linear constituents. The source reports verification for symmetric, alternating, \(p\)-solvable, and sporadic simple groups.
\subsection{Sources}
\begin{itemize}
\item Local UnsolvedMath registry, supplied complete source record: permanent ID \(701581\), exact source problem code \texttt{catalogue entry}, title \emph{Kourovka Notebook Problem 19.33}. The supplied status is \texttt{open}; this is a registry label, not an independently verified literature assessment.
\item \emph{The Kourovka Notebook (Unsolved Problems in Group Theory)}, Problem 19.33, proposed by E. Giannelli in 2018. The supplied background identifies the 19th issue (2018) and the 21st issue edition (2026), with identifier \texttt{arXiv:1401.0300v48}. The supplied collection description names E. I. Khukhro and V. D. Mazurov as editors.
\item The supplied statement reports verification for symmetric, alternating, \(p\)-solvable, and sporadic simple groups, without separate references for these verifications. No independent source consultation was performed.
\end{itemize}
\end{document}
